(i): $N_{W} = \sum_{i\le3} v_{i}w_{i} = \sum_{i\le3}(v_{i}\wedge w_{i} + \tfrac12)$ since $v_{i}\cdot w_{i} = \tfrac12$, so $1 - \tfrac23 N_{W} = -\tfrac23\sum v_{i}\wedge w_{i}$. For the rung: $w_{4}(1 - 2v_{4}w_{4}) = w_{4} - 2(1 - v_{4}w_{4})w_{4} = -w_{4}$ using $w_{4}^{2} = 0$, so $w_{4}\chi = -w_{4}(1 - 2n_{4})(1 - 2n_{5}) = w_{4}(1 - 2n_{5}) = -2w_{4}(v_{5}\wedge w_{5})$. The three factors $w_{4}$, $v_{5}\wedge w_{5}$, $\sum v_{i}\wedge w_{i}$ are built from mutually orthogonal vectors ($w_{4}$ anticommutes with $v_{5}, w_{5}$ and with every pair vector; the rung plane is orthogonal to the pair space), so their Clifford product is their wedge and has grade $1 + 2 + 2 = 5$. The script verifies the operator identity $T_{\mathrm{GJ}} = \tfrac43 w_{4}(v_{5}\wedge w_{5})(\sum v_{i}\wedge w_{i})$ on the module, the anticommutations, and, independently, the grade decomposition of $T_{\mathrm{GJ}}$ by the trace formula in the orthonormal basis $a_{i} = v_{i} + w_{i}$, $b_{i} = v_{i} - w_{i}$ of $\Cl(5,5)$ (coefficient of the blade $\Gamma_{A}$ is $\mathrm{tr}(\Gamma_{A}^{-1}T)/32$, the blades being trace-orthogonal in the irreducible representation), finding grade five only. The last clause of (i) is Fact \claimref{families.imported-decomposition-grade-dictionary} applied to the map $T \mapsto Y_{T}$, which is $\mathrm{Spin}(10)$-equivariant because $B$ is invariant: the grade-five space is $126 \oplus \overline{126}$, the bilinear forms on the sixteen are $16 \otimes 16 = 10 \oplus 120 \oplus 126$, and the $\overline{126}$ component, having no isomorphic summand among them, is sent to zero by Schur's lemma. (ii): with the species names that Hypothesis (S) supplies (\claimref{standard-model.hyp-s-slot-statistics}), $Y_{T_{\mathrm{GJ}}}(\psi, \varphi) = B(\psi, w_{4}\varphi)\cdot(B - L)(\varphi)\chi(\varphi)$ since $B - L$ and $\chi$ are diagonal; on $(L, e^{c})$ the factor is $(B-L)(e^{c})\chi(e^{c}) = 1\cdot(-1)$ in one order and $(B-L)(L)\chi(L) = (-1)(+1)$ in the other, on $(Q, d^{c})$ it is $(-\tfrac13)(-1)$ and $(\tfrac13)(+1)$: equal in both orders, which is the symmetry, with the stated values times the $\pm1$ of $Y_{w_{4}}$. The script enumerates the form and checks symmetry and hypercharge conservation entry by entry. (iii): without $\chi$ the two orders carry $(B-L)(e^{c}) = 1$ and $(B-L)(L) = -1$, opposite signs; the script checks antisymmetry and the neutrality sums. (iv) is the same computation with $v_{4}$; the script enumerates the form, checks its symmetry, and records the signed entries, $+1$ on both $(L, \nu^{c})$ entries and $-\tfrac13$, $+\tfrac13$ on the $(Q, u^{c})$ entries according to colour, which is why only the magnitude is stated.
